\section{Discussion}
\label{sec:discussion}
\subsection{An empirical diagnostic within an established problem}
Legend compatibility, heterogeneous annotations and class-space differences are established concerns \citep{yang2017harmonization,lambert2020mseg,kim2022uniseg}. Transferable land-cover semantics are also studied by TACOSS and LandSegmenter \citep{zermatten2025tacoss,liu2026landsegmenter}. The present contribution is a bounded empirical audit: connect the declared crosswalk to the masks actually available to a particular learner, then report the tested source mixture's recovery and supported-leaf cost under the specified controls. Neither taxonomy construction nor negative transfer is claimed as a new discovery.

Exact-label absence qualifies the available supervision; it does not exclude visual shift, establish a sole cause of failure or imply that a pretrained or heterogeneous-supervision model cannot recognize the concept. The primary CLIP and archived native U-Net comparisons show the same direction under their respective recorded training regimes. That extension narrows an explanation restricted to this one CLIP calibration instrument, but two chosen model families on a shared target do not establish architecture-wide generality.

\subsection{What source mixtures and pairing establish}
The primary experiment adds FLAIR imagery and herbaceous labels together. Exposure and duplicate-crop controls constrain two alternative explanations but do not turn that contrast into a label-only intervention. The recorded mismatched configuration retains most recovery, and native B1 gives a smaller additional contrast. Missing complete primary schedule files limit independent verification of pairing and exposure constraints. Auxiliary image content, regional statistics and weak alignment under same-region derangement remain compatible explanations. The retained recovery fraction is an arithmetic description, not a causal percentage attributable to a component. The retired U-Net and reversed-source shuffle arms provide no further efficacy claim.

A missing-leaf gain does not establish better aggregate transfer or preserved prior capabilities. Common/private performance reporting and known-class negative transfer already appear in adaptation research \citep{choe2024bus,choe2025unida,ren2026sats}. Here, the reported costs concern six exact LoveDA-supported leaves, or the three exact Semantic Drone-supported leaves in the reversed case. Decision-boundary changes, source appearance, sampling and loss weighting remain possible contributors. The experiments do not identify representational capacity as the cause and do not validate a mitigation module.

\subsection{Replication and provenance boundaries}
Semantic Drone supplies a second target replay using the primary CLIP checkpoints; its grass-versus-rest metric does not measure retention over all supported leaves. The reversed source-mixture case changes the missing leaf but reuses the OpenEarthMap target release. The native U-Net records change the architecture but retain the primary target population. These extensions are complementary checks with shared components, not a fully independent architecture-by-target replication. B2 establishes learnability under a different direct-supervision regime, not external-transfer superiority.

Semantic Drone's three-checkpoint points are reconstructable from archived integer totals, but per-raster records for two checkpoints were not recovered; its archived multi-checkpoint interval is not used. Native U-Net identifier sets and per-raster reference marginals match the primary population, but spatial label equality and original checkpoint execution are not thereby verified. The source-ready audit content is recovered by hash; this does not recover historical timestamp or executed-code identity. City/raster intervals condition on recorded checkpoints, not training-seed uncertainty or independent geographic sites. Geographic guards reduce named overlap risks without proving coordinate-level separation. Within these limits, a deployment report can make both the tested missing-leaf gain and its source-supported cost explicit.


\subsection{Class-proxy and construct-validity boundary}
The fixed codebook does not make the native annotation meanings identical. OEM rangeland includes bushes and shrubs, so its assignment to the herbaceous output is a coarse vegetation proxy (Supplementary Section S2). The reported rangeland recovery remains a native-label score under that declared mapping; it is not evidence of recovery restricted to pure herbaceous pixels. No grass/shrub subdivision is available in the retained OEM reference labels. Consequently, ontology resolution and annotation scope remain plausible contributors alongside imagery and supervision. We retain the original mappings, prompts, checkpoints and integer scores and do not synthesize a botanical reference or change model selection to resolve this limitation.
